\input{preamble}

\section*{Fermi's golden rule}

Find the transition rate $R_{a\rightarrow b}$ for the perturbing Hamiltonian
\begin{equation*}
H_1(\mathbf r,t)=2V(\mathbf r)\cos(\omega t+\phi)
\end{equation*}

\subsubsection*{\S1}

Let $\Psi(\mathbf r,t)$ be the following linear combination of two eigenstates.
\begin{equation*}
\Psi(\mathbf r,t)=c_a(t)\psi_a(\mathbf r)\exp(-i\omega_at)+c_b(t)\psi_b(\mathbf r)\exp(-i\omega_bt)
\end{equation*}

We need to solve for $c_b(t)$ to find the transition rate $R_{a\rightarrow b}$.

\bigskip

Start with the Schrodinger equation
\begin{equation*}
i\hbar\frac{\partial}{\partial t}\Psi(\mathbf r,t)=H_0(\mathbf r)\Psi(\mathbf r,t)+H_1(\mathbf r,t)\Psi(\mathbf r,t)
\end{equation*}

where
\begin{equation*}
H_0(\mathbf r)\Psi(\mathbf r,t)
=E_ac_a(t)\psi_a(\mathbf r)\exp(-i\omega_at)
+E_bc_b(t)\psi_b(\mathbf r)\exp(-i\omega_bt)
\end{equation*}

Evaluate the time derivative.
\begin{multline*}
i\hbar\frac{\partial}{\partial t}\Psi(\mathbf r,t)
=i\hbar\frac{\partial c_a(t)}{\partial t}\psi_a(\mathbf r)\exp(-i\omega_at)
+\boxed{\hbar\omega_ac_a(t)\psi_a(\mathbf r)\exp(-i\omega_at)}
\\
{}+i\hbar\frac{\partial c_b(t)}{\partial t}\psi_b(\mathbf r)\exp(-i\omega_bt)
+\boxed{\hbar\omega_bc_b(t)\psi_b(\mathbf r)\exp(-i\omega_bt)}
\\
{}=H_0(\mathbf r)\Psi(\mathbf r,t)+H_1(\mathbf r,t)\Psi(\mathbf r,t)
\tag{1}
\end{multline*}

We have $\hbar\omega_a=E_a$ and $\hbar\omega_b=E_b$
hence the boxed terms cancel with $H_0(\mathbf r)\Psi(\mathbf r,t)$ leaving
\begin{equation*}
i\hbar\frac{\partial c_a(t)}{\partial t}\psi_a(\mathbf r)\exp(-i\omega_at)
+i\hbar\frac{\partial c_b(t)}{\partial t}\psi_b(\mathbf r)\exp(-i\omega_bt)
=H_1(\mathbf r,t)\Psi(\mathbf r,t)
\end{equation*}

Substitute for $H_1(\mathbf r)$ and $\Psi(\mathbf r,t)$.
\begin{multline*}
i\hbar\frac{\partial c_a(t)}{\partial t}\psi_a(\mathbf r)\exp(-i\omega_at)
+i\hbar\frac{\partial c_b(t)}{\partial t}\psi_b(\mathbf r)\exp(-i\omega_bt)
\\
{}=2V(\mathbf r)\cos(\omega t+\phi)
\biggl[
c_a(t)\psi_a(\mathbf r)\exp(-i\omega_at)+c_b(t)\psi_b(\mathbf r)\exp(-i\omega_bt)
\biggr]
\end{multline*}

Apply an inner product with $\psi_b^*(\mathbf r)$ to eliminate wave functions.
\begin{equation*}
i\hbar\frac{\partial c_b(t)}{\partial t}\exp(-i\omega_bt)
=2\cos(\omega t+\phi)
\biggl[
c_a(t)M_{21}\exp(-i\omega_at)+
c_b(t)M_{22}\exp(-i\omega_bt)
\biggr]
\end{equation*}

Symbols $M_{21}$ and $M_{22}$ are the matrix elements
\begin{align*}
M_{21}&=\int\psi_b^*(\mathbf r)V(\mathbf r)\psi_a(\mathbf r)\,d\mathbf r
\\
M_{22}&=\int\psi_b^*(\mathbf r)V(\mathbf r)\psi_b(\mathbf r)\,d\mathbf r
\end{align*}

Multiply both sides by $\exp(i\omega_bt)$.
\begin{equation*}
i\hbar\frac{\partial c_b(t)}{\partial t}
=2\cos(\omega t+\phi)
\biggl[
c_a(t)M_{21}\exp[i(\omega_b-\omega_a)t]+c_b(t)M_{22}
\biggr]
\end{equation*}

For time $t$ near the origin use $c_a(t)\approx1$ and $c_b(t)\approx0$ to obtain
\begin{equation*}
i\hbar\frac{\partial c_b(t)}{\partial t}
=2\cos(\omega t+\phi)M_{21}\exp[i(\omega_b-\omega_a)t]
\end{equation*}

Solve for $c_b(t)$ by integrating.
\begin{equation*}
c_b(t)=\frac{2M_{21}}{i\hbar}
\int_0^t\cos(\omega t'+\phi)\exp[i(\omega_b-\omega_a)t']\,dt'
\end{equation*}

In exponential form
\begin{multline*}
c_b(t)=\frac{M_{21}}{i\hbar}
\int_0^t
\exp(i\omega t'+\phi)\exp[i(\omega_b-\omega_a)t']\,dt'
\\
{}+\frac{M_{21}}{i\hbar}
\int_0^t
\exp(-i\omega t'-\phi)\exp[i(\omega_b-\omega_a)t']\,dt'
\end{multline*}

The solution is
\begin{multline*}
c_b(t)
=-\frac{M_{21}}{\hbar}
\,
\frac{\exp[i(\omega_b-\omega_a+\omega) t]-1}{\omega_b-\omega_a+\omega}
\,
\exp(i\phi)
\\{}-\frac{M_{21}}{\hbar}
\,
\frac{\exp[i(\omega_b-\omega_a-\omega) t]-1}{\omega_b-\omega_a-\omega}
\,
\exp(-i\phi)
\tag{2}
\end{multline*}

For $\omega$ such that $\omega\approx\omega_b-\omega_a$ the second term
dominates so discard the first term.
\begin{equation*}
c_b(t)=-\frac{M_{21}}{\hbar}
\,
\frac{\exp[i(\omega_b-\omega_a-\omega) t]-1}{\omega_b-\omega_a-\omega}
\,
\exp(-i\phi)
\end{equation*}

By the identity
\begin{equation*}
\exp(ix)-1=2i\sin\biggl(\frac{x}{2}\biggr)\exp\biggl(\frac{ix}{2}\biggr)
\tag{3}
\end{equation*}

we have
\begin{equation*}
c_b(t)=-\frac{2iM_{21}}{\hbar}
\,
\frac{
\sin\left[\frac{1}{2}(\omega_b-\omega_a-\omega)t\right]
\exp\left[\frac{1}{2}i(\omega_b-\omega_a-\omega)t\right]
}{\omega_b-\omega_a-\omega}
\,
\exp(-i\phi)
\end{equation*}

Multiply numerator and denominator by $t/2$ to prepare for sinc function.
\begin{equation*}
c_b(t)=-\frac{itM_{21}}{\hbar}
\,
\frac{
\sin\left[\frac{1}{2}(\omega_b-\omega_a-\omega)t\right]
\exp\left[\frac{1}{2}i(\omega_b-\omega_a-\omega)t\right]
}{\frac{1}{2}(\omega_b-\omega_a-\omega)t}
\,
\exp(-i\phi)
\end{equation*}

Substitute $\operatorname{sinc}(x)$ for $\sin(x)/x$.
\begin{equation*}
c_b(t)=-\frac{itM_{21}}{\hbar}
\operatorname{sinc}\left[\frac{1}{2}(\omega_b-\omega_a-\omega)t\right]
\exp\left[\frac{1}{2}i(\omega_b-\omega_a-\omega)t\right]
\exp(-i\phi)
\tag{4}
\end{equation*}

\subsubsection*{\S2}

Transition density $P$ is the squared magnitude of $c_b(t)$.
\begin{equation*}
P=|c_b(t)|^2=\frac{t^2}{\hbar^2}\,|M_{21}|^2
\operatorname{sinc}^2\left[\frac{1}{2}(\omega_b-\omega_a-\omega)t\right]
\tag{5}
\end{equation*}

Total density $P_{tot}$ is the integral of $P$ over a small energy band.
\begin{equation*}
P_{tot}=\frac{t^2}{\hbar^2}|M_{21}|^2
\int_{E-\epsilon}^{E+\epsilon}
\operatorname{sinc}^2\left[\frac{1}{2}(E'/\hbar-\omega)t\right]
\rho(E')\,dE'
\end{equation*}

The center of the energy band is $E=\hbar(\omega_b-\omega_a)$.
Symbol $\rho$ is density of states.

\bigskip

Use the approximation $\rho(E')\approx\rho(\hbar\omega)$ to obtain
\begin{equation*}
P_{tot}=\frac{t^2}{\hbar^2}|M_{21}|^2\rho(\hbar\omega)
\int_{E-\epsilon}^{E+\epsilon}
\operatorname{sinc}^2\left[\frac{1}{2}(E'/\hbar-\omega)t\right]\,dE'
\end{equation*}

Let $y$ be a change of variable such that
\begin{equation*}
y=\frac{1}{2}(E'/\hbar-\omega)t
\end{equation*}

From the definition of $y$ we have
\begin{equation*}
E'=\frac{2\hbar y}{t}+\hbar\omega
\end{equation*}

hence
\begin{equation*}
dE'=\frac{2\hbar}{t}\,dy
\end{equation*}

The integral limits transform as
\begin{equation*}
E\pm\epsilon
\rightarrow
\frac{1}{2}[(E\pm\epsilon)/\hbar-\omega]t
=\frac{Et}{2\hbar}-\frac{\omega t}{2}\pm\frac{\epsilon t}{2\hbar}
=\pm\frac{\epsilon t}{2\hbar}
\end{equation*}

After change of variable from $E'$ to $y$ we have
(note $t^2/\hbar^2$ is reduced to $t/\hbar$)
\begin{equation*}
P_{tot}=\frac{2t}{\hbar}|M_{21}|^2\rho(\hbar\omega)
\int_{-\epsilon t/2\hbar}^{\epsilon t/2\hbar}
\operatorname{sinc}^2y\,dy
\end{equation*}

Use the approximation
\begin{equation*}
\int_{-\epsilon t/2\hbar}^{\epsilon t/2\hbar}\operatorname{sinc}^2 y\,dy
\approx
\int_{-\infty}^\infty\operatorname{sinc}^2 y\,dy=\pi
\end{equation*}

to obtain
\begin{equation*}
P_{tot}=\frac{2\pi t}{\hbar}|M_{21}|^2\rho(\hbar\omega)
\end{equation*}

The transition rate is the derivative of $P_{tot}$.
\begin{equation*}
R_{a\rightarrow b}
=\frac{d}{dt}P_{tot}=\frac{2\pi}{\hbar}|M_{21}|^2\rho(\hbar\omega)
\tag{6}
\end{equation*}

\subsubsection*{\S3}

We started with
\begin{equation*}
H_1(\mathbf r,t)=2V(\mathbf r)\cos(\omega t+\phi)
\end{equation*}

and obtained Fermi's golden rule (6).
Consequently for $H_1(\mathbf r,t)$ such that
\begin{equation*}
H_1(\mathbf r,t)=V(\mathbf r)\cos(\omega t+\phi)
\end{equation*}

we would have
\begin{equation*}
R_{a\rightarrow b}=\frac{2\pi}{\hbar}\left|\frac{M_{21}}{2}\right|^2\rho(\hbar\omega)
\end{equation*}

\subsubsection*{\S4}

\href{https://georgeweigt.github.io/examples/fermi-golden-rule-demo.html}{Eigenmath script}

\end{document}
